\begin{table*}[!t]\centering\scriptsize
\caption{Complementary validation axes in the formal design.}\label{tab:validation}
\begin{tabularx}{.98\textwidth}{L{3.10cm}L{2.45cm}Y Y}\toprule
Evidence axis & Grouping unit & Fit/evaluation design & Statistical role \\\midrule
Spatial GroupKFold & 5 km block & Separate folds within historical diagnostic samples; adjacent blocks can occur on opposite fold sides & Block-grouped within-domain spatial diagnostic, not buffered geographic holdout \\
LOYO & Constituent year & Fit on other window years; evaluate the temporally unseen year & Internal temporal cross-validation; no window ranking \\
Historical spatial reserve & 5 km block & Fit on diagnostic blocks; evaluate once on whole blocks excluded from GroupKFold and LOYO & Pre-refit held-out historical spatial diagnostic; not a final-model test \\
Multi-AOI target & AOI $\times$ sensor & Refit each of six historical windows; apply to 2025 & Cross-domain replication on the common target grid; effective valid support can differ by sensor \\
Rolling-Origin target & Target year & Refit H1--H3 using only preceding years; apply to 2024 or 2025 & Chronological forward-transfer evaluation \\
2025 DPM benchmark & AOI $\times$ sensor & Four prespecified empirical NDVI quantile pairs define lower/upper endmember proxies; target-year NDVI, but no target-year FCOVER numeric values, enters endpoint estimation; the minimum-RMSE configuration is shown retrospectively & Secondary descriptive benchmark under a distinct information condition \\
Block inference & 5 km block & Three paired H1--H3 contrasts per sensor--AOI--target family on identical target identities within each paired block; Holm correction within each family & Local historical-window inference; cross-family totals descriptive \\\bottomrule
\end{tabularx}
\end{table*}
